\section{Conclusions}
\label{sec:conclusion}

We analyzed 21~million GPS-instrumented e-scooter trips from 52 South Korean cities using a predict-then-validate framework to evaluate how software speed governance reshapes riding behavior beyond the tautological speed reduction. Cross-sectional behavioral profiling of eight genuinely behavioral outcomes across safety, demand, and riding dynamics generated six testable predictions about which margins would respond to governance and which would not. The December 2023 system-wide removal of the ungoverned mode provided a natural experiment to validate each prediction via multi-outcome difference-in-differences, within-user compensation testing, and dose-response analysis.

Three principal findings emerge. First, cross-sectional behavioral profiling generates accurate predictions about causal governance effects: all six predictions---derived from observational mode effect sizes, within-subject paired comparisons, and SHAP feature importance---are confirmed by the natural experiment. Outcomes with strong cross-sectional mode dependence (harsh acceleration $\beta = -0.61$, harsh deceleration $\beta = -0.63$) respond causally (Bonferroni-corrected, dose-response across 50 cities), while outcomes with weak or no mode dependence (speed CV, cruise fraction, trip counts, user retention) show null causal effects. Second, there is zero behavioral compensation on any unconstrained margin: the within-user mode-switcher analysis (24,357 switchers vs.\ 79,700 controls) yields negligible effect sizes on all six outcomes ($|d| < 0.14$), confirming Prediction P6. Third, the predict-then-validate framework provides a methodological template for pre-implementation assessment of speed management policies: the cross-sectional mode effect size serves as a reliable leading indicator for causal policy response.

The behavioral escalation pathway analysis reveals that 60--72\% of the experience-related increase in harsh events is mediated by progressive TUB adoption, not by within-mode behavioral change, suggesting that graduated mode access could interrupt the escalation pathway. The predict-then-validate concordance has a reassuring policy implication: cities considering governor mandates can use observational behavioral profiling to forecast effects before implementation, and the forecast is that governor mandates are Pareto-improving on measured margins.

Future work should extend these findings in four directions: (1)~linking speed governance to crash outcomes using hospital or insurance records, (2)~replicating the predict-then-validate framework with operators in other countries to test generalizability, (3)~evaluating dynamic geofencing as a context-sensitive alternative to uniform speed caps, and (4)~using higher-frequency kinematic data (accelerometers) to capture a broader range of aggressive riding behaviors.
